% ═══════════════════════════════════════════════════════════════════════
%  PREÀMBUL COMPARTIT DEL BANC DE PREGUNTES
%  ─────────────────────────────────────────────────────────────────────
%  CONTRACTE (llegeix-ho abans de tocar res):
%
%  · Aquest fitxer NO es compila sol. És un fragment.
%  · Qui el fa servir (build.py o el lloc web) SEMPRE emet, per aquest
%    ordre:  \documentclass → \newif\ifsolucions + \solucions(true|false)
%            → \newif\ifcurt + \curt(true|false) → aquest fitxer
%            → \begin{document} → cossos → \end{document}
%  · \ifcurt diu la modalitat: examen d'1 h 30 (false) o de 50 min (true).
%  · Tota pregunta del banc compila amb AQUEST preàmbul i cap altre.
%    Si una pregunta necessita un paquet nou, s'afegeix aquí i es torna
%    a compilar TOT el banc. Mai un \usepackage dins d'una pregunta.
% ═══════════════════════════════════════════════════════════════════════

\usepackage[T1]{fontenc}
\usepackage{lmodern}
\usepackage[catalan]{babel}
\usepackage{amsmath,amssymb}
\usepackage{array}
\usepackage{tikz}
\usepackage{enumitem}
\usepackage{xcolor}
\usepackage{comment}
\usepackage{needspace}
\usepackage[a4paper,top=1.8cm,bottom=1.8cm,left=2cm,right=2cm]{geometry}

\setlength{\parindent}{0pt}
\setlength{\parskip}{3pt}

% ── 1. Capçalera de pregunta ──────────────────────────────────────────
% El filet va A SOBRE de cada pregunta i el \Needspace el manté enganxat
% al seu text: així el filet mai queda orfe al capdamunt d'una pàgina.
\newcommand{\encapcalament}[1]{%
  \par\Needspace{8\baselineskip}%
  \bigskip\hrule\medskip
  \noindent\textbf{#1}\par\smallskip}

% ── 2. Apartats a) b) c) amb la seva puntuació ────────────────────────
% \apartat{0,75} escriu "(0,75 punts)".  build.py llegeix aquests valors
% del .tex: la puntuació viu AQUÍ i enlloc més. L'opcional, si hi és, és la
% puntuació a l'examen de 50 min: \apartat[1,25]{0,75}.
\makeatletter
\newcommand{\bp@ptsmot}[1]{\def\bp@a{#1}\def\bp@u{1}%
  \ifx\bp@a\bp@u 1~punt\else #1~punts\fi}
% Modalitat: \apartat[1,25]{0,75} val 0,75 punts a l'examen d'1 h 30 i 1,25 al
% de 50 min. Sense l'opcional, la mateixa puntuació a tots dos.
\newcommand{\apartat}[2][]{%
  \ifcurt\if\relax\detokenize{#1}\relax\def\bp@v{#2}\else\def\bp@v{#1}\fi
  \else\def\bp@v{#2}\fi
  \item \textit{(\expandafter\bp@ptsmot\expandafter{\bp@v})}\enspace\ignorespaces}
\makeatother
\newenvironment{apartats}
  {\begin{enumerate}[label=\textbf{\alph*)},leftmargin=*,itemsep=\bigskipamount,
                     topsep=\medskipamount,parsep=\smallskipamount]}
  {\end{enumerate}}

% ── 3. Graella de subapartats i) ii) iii) en mode display ─────────────
% \begin{graella}{4} \sa ... & \sa ... \\ \sa ... \end{graella}
% La columna força \displaystyle: els \lim hi surten amb el subíndex A
% SOTA, igual que en una fórmula destacada. Mai fem servir array pelat.
\newcounter{bpsa}
\newcommand{\sa}{\stepcounter{bpsa}\textup{\roman{bpsa})}~}
\newenvironment{graella}[1]
  {\setcounter{bpsa}{0}\setlength{\arraycolsep}{1.2em}%
   \[\begin{array}{*{#1}{>{\displaystyle}l}}}
  {\end{array}\]}

% ── 3b. Fórmules destacades: sempre el mateix aire ───────────────────
% Si la línia d'abans és curta («determina:»), TeX fa servir un espai
% «curt», gairebé zero, i la fórmula queda enganxada al text. L'igualem
% a l'espai normal.
\makeatletter
\g@addto@macro\normalsize{%
  \setlength\abovedisplayshortskip{\abovedisplayskip}%
  \setlength\belowdisplayshortskip{\belowdisplayskip}}
\makeatother
\AtBeginDocument{\normalsize}

% ── 4. Condicions dins de \begin{cases} ───────────────────────────────
% \si{-1\le x\le 2}  →  el signe menys surt unari i ben espaiat.
% Escriure \text{si } -1\le x\le 2 a pèl produeix "si − 1 ≤ x ≤ 2". Bug.
\newcommand{\si}[1]{\text{si }{#1}}

% ── 5. Solucions ──────────────────────────────────────────────────────
% REGLA: \end{solucio} ha d'anar SOL a la seva línia (ho exigeix el
% paquet comment quan la solució s'exclou). build.py ho comprova.
\ifsolucions
  \newenvironment{solucio}
    {\par\nopagebreak\medskip\begingroup\color{blue!55!black}\small
     \textbf{Solució.}\enspace\ignorespaces}
    {\par\endgroup\smallskip}
\else
  \excludecomment{solucio}
\fi

% ── 5b. Apartats que només surten a l'examen d'1 h 30 ─────────────────
% \begin{nomesllarg} … \end{nomesllarg} envolta un apartat sencer, amb la
% seva solució. A l'examen de 50 min no hi surt, i els apartats de després
% es tornen a lletrejar sols. \end{nomesllarg} va SOL a la seva línia.
\ifcurt \excludecomment{nomesllarg} \else \includecomment{nomesllarg} \fi

% ── 6. Procedència de les preguntes PAU ───────────────────────────────
% La línia «PAU juny 2026, sèrie 1» la injecta el build (i el lloc) just
% després de la capçalera, a partir del codi de la pregunta i de
% pau/convocatories.json. No s'escriu mai a mà: build.py ho rebutja.
\newcommand{\procedencia}[1]{\noindent{\small\itshape #1}\par\vspace{3pt}}
